\section{What the Paths Mean: Real Indications on the Full Graph}
\label{sec:indications}

The benchmark of Sections~\ref{sec:blind}--\ref{sec:targeted} asks whether a
typed path survives. It does not ask whether the path says anything true, and
it lives on a BFS sample. This section removes both limitations. We use
PrimeKG's own list of approved drug indications as ground truth and run the
analysis on the full filtered graph.

\subsection{Design}
\label{sec:indications-design}

PrimeKG lists 9,388 drug--disease indication pairs. For 6,233 of them, both the
drug and the disease are in the filtered graph. These pairs were never part of
the graph we analyze, because Section~\ref{sec:data-graphs} kept only
drug--protein, disease--protein and protein--protein edges. They are therefore
external evidence that the retrieval rule has not seen. For each indication we
form a decoy by permuting the diseases among the indicated drugs and rejecting
any pair that PrimeKG lists as an indication, a contraindication or an
off-label use. Every decoy keeps the drug of one indication, and the decoy diseases
are the indicated diseases in a new order, so the two sets have almost
identical distributions of drug degree and disease degree. They differ in
whether the pair is a recorded therapeutic relation. Retrieval at hop limit $L$ is defined as in
Section~\ref{sec:data-benchmark}. We report the recall (the share of
indications with a typed path), the decoy rate (the same share for decoys),
their ratio, and the mutual information in bits between the label and the
retrieval outcome for a balanced mixture of the two sets. Damage is random edge
removal (five orders) or removal of nodes in decreasing order of degree, both on
the full graph of 30,968 nodes.

\subsection{Only short paths are specific}
\label{sec:indications-specific}

Table~\ref{tab:indications} and Figure~\ref{fig:summary}b show the intact
graph. A shared protein links 31.0\% of indications and 14.2\% of decoys, so a
two-hop answer is 2.2 times more common among indications than among decoys. The
contrast fades quickly with the horizon. At $L=3$ the ratio is 1.26, at $L=4$
it is 1.01, and without a limit both sets are connected in 99.97\% of cases. A
path of unbounded length carries no information about indications at all. These
ratios are stable. Over five further decoy draws the two-hop ratio lies between
2.09 and 2.18, and a bootstrap over the 1,375 indicated drugs, which accounts for
indications that share a drug, gives 95\% intervals of 1.92--2.47 at $L=2$,
1.20--1.32 at $L=3$ and 1.00--1.02 at $L=4$.

\begin{table}[!htbp]
\centering
\small
\caption{Typed retrieval of 6,233 approved indications and 6,233
degree-preserving decoys on the intact full graph. The configuration-model
column is the expected decoy rate in a random graph with the same layer degrees
(Eq.~\eqref{eq:cm}). Information is per pair, for a balanced mixture.}
\label{tab:indications}
\begin{tabular}{@{}lrrrrr@{}}
\toprule
Hop limit & Indications & Decoys & Ratio & Configuration model & Information (mbit) \\
\midrule
$L=2$      & 31.0\% & 14.2\% & 2.18 & 13.1\% & 29.6 \\
$L=3$      & 61.1\% & 48.3\% & 1.26 & 58.1\% & 11.9 \\
$L=4$      & 96.2\% & 94.9\% & 1.01 & 99.7\% & 0.7 \\
Unlimited  & 99.97\% & 99.97\% & 1.00 & -- & 0.0 \\
\bottomrule
\end{tabular}
\end{table}

The loss of specificity is the small-world property, and the generating-function
picture of Section~\ref{sec:theory-gf} predicts it. In a configuration model
with the same layer degrees, a drug with $k_d$ targets and a disease with $k_y$
proteins share on average
\begin{equation}
\Lambda_2 = k_d k_y\, c_2,\qquad c_2=\frac{\sum_x k_D(x)\,k_Y(x)}{K_D K_Y}
\label{eq:cm}
\end{equation}
proteins, where $k_D(x)$ and $k_Y(x)$ are the drug and disease degrees of protein
$x$, as in Eq.~\eqref{eq:mf}, and $K_D$ and $K_Y$ are the layer totals. Paths
with one protein--protein step add $k_d k_y c_3$ with
$c_3=\sum_x k_D(x)k_x\sum_x k_Y(x)k_x/(K_D K_Y \sum_x k_x)$, where $k_x$ is the
PPI degree, and each further step multiplies the last term by $\kappa-1$, with
$\kappa=\langle k^2\rangle/\langle k\rangle=187$ for the PPI layer, because a
path enters a protein through one of its edges and can leave through any of the
others (Section~\ref{sec:theory-gf}). Thus $c_3/c_2=24$ and the four-hop term is
186 times the three-hop one. With $\Lambda_L$ the cumulative expected number of
paths of length at most $L$, the probability that at least one exists is
$1-e^{-\Lambda_L}$. This estimate,
which has no free parameters, gives a decoy rate of 13.1\% at $L=2$ against the
observed 14.2\%. It overestimates $L=3$ (58\% against 48\%), plausibly because it ignores
degree correlations, and it gives near-certain connection at $L=4$. At $L=2$,
decoys therefore behave like random pairs, while indications are 2.2 times more
likely than chance to share a protein. The excess is real biology, but the chance
paths of longer horizons swamp it.

Contraindications sit in between. A shared protein links 20.2\% of the 21,835
contraindicated pairs in the graph. A two-hop path therefore signals a
mechanistic relation between a drug and a disease, not necessarily a
therapeutic one.

\subsection{The specific paths are the fragile ones}
\label{sec:indications-fragile}

The information sits in exactly the observable that Sections~\ref{sec:theory}
and~\ref{sec:blind} found to be fragile. Under random edge removal
(Figures~\ref{fig:summary}a and~\ref{fig:indications}a), the giant component of
the full graph still holds 86\% of nodes at $f=0.5$, and 85\% of all
indications remain connected without a hop limit. Two-hop recall falls from
31.0\% to 12.0\%. Of the 1,931 indications that had a two-hop link, 40.3\% keep
it (mean of 500 orders), as Eq.~\eqref{eq:r2} predicts from their
shared-protein counts (0.403); on the sample benchmark the corresponding value
was 0.41. A single damage state scatters around this mean with a standard
deviation of 0.020, twice the independent-query value, because indications
share drugs. The ratio of
indication recall to decoy rate rises only slowly, from 2.2 to 2.4 at $f=0.5$,
but fewer pairs have any answer, so the information carried by a two-hop answer
falls to 40\% of its intact value (Figure~\ref{fig:indications}b). This decay
is predicted without parameters. Applying Eq.~\eqref{eq:r2} separately to the
shared-protein counts of indications and of decoys gives both rates, and hence
the information, at every damage level; the prediction stays within 1 millibit
of the simulation, inside the spread of a single damage state. For $L=3$ the
independent-path estimate of Section~\ref{sec:blind-horizon}, applied to each
pair's actual counts of two- and three-hop paths, does almost as well (within
1.4 millibits). Longer
horizons behave differently. The three-hop ratio rises from 1.26 to 2.0 at
$f=0.8$, where a three-hop answer carries as much information as a two-hop one
(3.9 against 3.5 millibits, within one standard deviation). Thinning reduces the
number of chance paths below one per pair, and the presence of a path again
reflects how many paths a pair has. The advantage of two-hop answers over
three-hop answers therefore disappears by $f\approx0.8$.

\begin{figure}[!htbp]
  \centering
  \includegraphics[width=\linewidth]{act6/fig_indications.pdf}
  \caption{All 6,233 approved indications and their decoys on the full filtered
  graph under random edge removal. (a) Giant component, and recall of
  indications (solid) and decoys (dotted) at $L=2$, 3 and unlimited; the dashed line is
  Eq.~\eqref{eq:r2} for indications and for decoys. Means of five damage orders.
  (b) Mutual information between the label and the retrieval outcome, in
  millibits per pair. Shading shows one standard deviation over orders. Black
  lines are predictions with no fitted parameters: Eq.~\eqref{eq:r2} for $L=2$
  and the independent-path estimate for $L=3$.}
  \label{fig:indications}
\end{figure}

\subsection{Hubs carry the non-specific paths}
\label{sec:indications-hubs}

Targeted damage tells the same story from the other side
(Figure~\ref{fig:indications-attack}). Of the shared proteins that link decoys,
45\% are among the 1\% of proteins with the highest degree (degree at least
316). For indications the share is 24\%. Removing the top 1\% of nodes by degree
cuts two-hop recall from 31.0\% to 20.0\%, but it cuts the decoy rate from
14.2\% to 5.0\%. The likelihood ratio rises from 2.2 to 4.0, and the information
per pair increases from 29.6 to 39.4 millibits. After removing 10\% of nodes,
only 2.5\% of indications keep a two-hop answer, but the ratio reaches 9.8.
Only 16 decoys keep a two-hop path at that point, so the ratio
becomes noisy and we do not interpret larger attacks.
The size-biased mean field of Eq.~\eqref{eq:mf} shows the same thing. After
removing 5\% of nodes it predicts route survival for the two-hop indications
about as well as for the sample (0.283 against 0.305 measured). After removing
1\%, however, it errs in the opposite direction from the sample: 78\% of the
two-hop indications whose endpoints survive keep a route, against 65\%
predicted, because only 27\% of their routes pass through the removed nodes,
against 44\% of all drug--disease routes in the graph. Real indications avoid
hubs. Hubs therefore provide much of the redundancy that keeps short paths alive
under blind damage, and many of the paths that connect unrelated pairs. Network
medicine corrects for this by comparing a drug--disease proximity with
degree-preserving random references \cite{guney2016network}. Our decoys play
the same role here.

\begin{figure}[!htbp]
  \centering
  \includegraphics[width=0.86\linewidth]{act6/fig_indications_attack.pdf}
  \caption{Degree attack on the full filtered graph. (a) Giant component, and
  recall of indications (solid) and decoys (dotted) at $L=2,3,4$. (b) Ratio of
  indication recall to decoy rate, on a logarithmic axis, up to 10\% of nodes
  removed.}
  \label{fig:indications-attack}
\end{figure}

\subsection{What this adds}
\label{sec:indications-adds}

The full graph and an external ground truth confirm the central finding. Global
connectivity barely moves while the short paths that answer queries disappear,
and Eq.~\eqref{eq:r2} predicts how fast they disappear. The experiment also
explains why this matters. Unbounded reachability is robust because almost
every pair is joined by many chance paths through hubs, and for the same reason
it says nothing about which drug treats which disease. The specific quantity, a
short typed path, is fragile because most two-hop indications are supported by
one or two shared proteins. Robustness
and specificity trade against each other in this graph, and a robustness
statistic is only meaningful for the observable that carries the information.
